\begin{frame}[t]{\ev{\tagO}Fork saves resources when preparation dominates reuse cost}
\vspace{.1cm}
\begin{center}
{\Large$(N-1)\,P \;>\; H + N\,(R+D)$\par}
\vspace{.15cm}
{\small preparation $P$, capture $H$, restore $R$, divergence $D$, children $N$. Additive resource-cost model.\par}
\end{center}
\vspace{.1cm}
{\small\color{Muted}Calculation, not a measurement. Units: resource-seconds.\par}
\vspace{.05cm}
{\small
\begin{tabularx}{\textwidth}{@{}L{4.6cm}L{3.0cm}L{2.9cm}Y@{}}
\toprule
Case & Avoided preparation & Added reuse cost & Result\\
\midrule
$N{=}8$, $P{=}30$, $H{=}2$, $R{=}1$, $D{=}3$ & $7\times30 = 210$ & $2 + 8\times4 = 34$ & fork saves 176\\
same, warm cache $P{=}2$ & $7\times2 = 14$ & $34$ & fork loses 20\\
\bottomrule
\end{tabularx}\par}
\vspace{.15cm}
{\small Batch elapsed time is a separate question: it depends on critical paths and contention. A faster-booting minimal image is the alternative (Jitsu, LightVM).\par}
\sources{Madhavapeddy et al.\ NSDI'15 (Jitsu); Manco et al.\ SOSP'17 (LightVM, NEC Labs). Parameter values are illustrative.}
\note{[0:45] Fork avoids N minus one preparations. It pays one capture, plus a restore and a divergence for each child. With eight children and a thirty-second preparation, fork avoids 210 resource-seconds and adds 34. It saves 176. With a warm cache, preparation falls to two seconds and fork loses. These are calculations. Elapsed time depends on critical paths and contention. Jitsu and LightVM show the faster-boot alternative. Next, we show what our measurements cover.}
\end{frame}
